\textbf{Phần 1. Trắc nghiệm lựa chọn}
\setcounterpageref{ex}{0}

\begin{ex} %[3D1V49]
Số liền sau của số lớn nhất có bốn chữ số là số nào?
\choice
{$9998$}
{$10\,001$}
{\True $10\,000$}
{$9990$}
\loigiai{
Số lớn nhất có 4 chữ số là 9999. Số liền sau của nó là $9999 + 1 = 10\,000$.
\textbf{Đáp án của câu hỏi là C}
}
\end{ex}

\begin{ex} %[3D1V49]
Trong các số sau, số nào khi làm tròn đến hàng trăm thì được số $4500$?
\choice
{\True $4485$}
{$4555$}
{$4449$}
{$4560$}
\loigiai{
Số 4485 có hàng chục là 8 ($8 \ge 5$) nên làm tròn lên $4500$.
\textbf{Đáp án của câu hỏi là A}
}
\end{ex}

\begin{ex} %[3D1V49]
Số La Mã $\text{XIX}$ cộng thêm $\text{II}$ được kết quả là số La Mã nào?
\choice
{$\text{XX}$}
{$\text{XXII}$}
{$\text{XVIII}$}
{\True $\text{XXI}$}
\loigiai{
$\text{XIX} = 19$; $\text{II} = 2$. $19 + 2 = 21 = \text{XXI}$.
\textbf{Đáp án của câu hỏi là D}
}
\end{ex}

\begin{ex} %[3D1V49]
Tìm $x$ là số tròn nghìn thỏa mãn: $6500 < x < 8500$. Giá trị của $x$ là
\choice
{$7000$}
{\True $7000$ hoặc $8000$}
{$8000$}
{$7500$}
\loigiai{
Các số tròn nghìn trong khoảng này là 7000 và 8000.
\textbf{Đáp án của câu hỏi là B}
}
\end{ex}

\begin{ex} %[3D1V49]
Từ các chữ số $3, 0, 7, 5$, lập được số có bốn chữ số khác nhau nhỏ nhất chia hết cho 5 là
\choice
{$3057$}
{$3570$}
{\True $3075$}
{$3750$}
\loigiai{
Hàng nghìn nhỏ nhất khác 0 là 3, hàng trăm 0. Để chia hết cho 5 thì tận cùng là 5 (nhỏ hơn 0 thì không thể vì 0 đã ở hàng trăm), hàng chục là 7. Số đó là $3075$.
\textbf{Đáp án của câu hỏi là C}
}
\end{ex}

\begin{ex} %[3D1V49]
Số gồm 8 nghìn, 5 chục và 3 đơn vị được viết là
\choice
{\True $8053$}
{$8530$}
{$8503$}
{$8035$}
\loigiai{
Chữ số hàng trăm bằng 0. Số đó là $8053$.
\textbf{Đáp án của câu hỏi là A}
}
\end{ex}

\begin{ex} %[3D1V49]
Làm tròn số $7249$ đến hàng chục được số $A$, làm tròn số $7249$ đến hàng trăm được số $B$. Giá trị của $A - B$ là
\choice
{$40$}
{$60$}
{$0$}
{\True $50$}
\loigiai{
$A = 7250$ (đơn vị 9 làm tròn lên). $B = 7200$ (hàng chục 4 làm tròn xuống). $A - B = 7250 - 7200 = 50$.
\textbf{Đáp án của câu hỏi là D}
}
\end{ex}

\begin{ex} %[3D1V49]
Một đồng hồ cổ có các số ghi bằng chữ số La Mã. Kim dài chỉ số $\text{VI}$, kim ngắn ở chính giữa số $\text{IV}$ và số $\text{V}$. Lúc đó là mấy giờ?
\choice
{$5$ giờ $30$ phút}
{\True $4$ giờ $30$ phút}
{$6$ giờ $20$ phút}
{$4$ giờ $6$ phút}
\loigiai{
Kim dài chỉ $\text{VI}$ là 30 phút. Kim ngắn giữa $\text{IV}$ (4) và $\text{V}$ (5) là 4 giờ 30 phút.
\textbf{Đáp án của câu hỏi là B}
}
\end{ex}

\begin{ex} %[3D1V49]
Số chẵn lớn nhất có bốn chữ số khác nhau là
\choice
{$9998$}
{$9874$}
{\True $9876$}
{$9864$}
\loigiai{
Hàng nghìn 9, trăm 8, chục 7, đơn vị là chữ số chẵn khác nhau lớn nhất là 6: $9876$.
\textbf{Đáp án của câu hỏi là C}
}
\end{ex}

\begin{ex} %[3D1V49]
Có bao nhiêu số có bốn chữ số lớn hơn $9990$?
\choice
{\True $9$ số}
{$10$ số}
{$8$ số}
{$11$ số}
\loigiai{
Các số có 4 chữ số lớn hơn 9990 là từ 9991 đến 9999. Có $9999 - 9991 + 1 = 9$ số.
\textbf{Đáp án của câu hỏi là A}
}
\end{ex}

\begin{ex} %[3D1V49]
Một cửa hàng bán được $4350\text{ kg}$ gạo trong tháng Một. Bác chủ cửa hàng làm tròn số gạo đến hàng trăm để báo cáo. Số gạo báo cáo là
\choice
{$4300\text{ kg}$}
{$4350\text{ kg}$}
{$5000\text{ kg}$}
{\True $4400\text{ kg}$}
\loigiai{
Chữ số hàng chục là 5 ($5 \ge 5$) nên làm tròn lên $4400\text{ kg}$.
\textbf{Đáp án của câu hỏi là D}
}
\end{ex}

\begin{ex} %[3D1V49]
Sắp xếp theo thứ tự giảm dần: $\text{XIX}, \text{XV}, \text{XXI}, \text{IX}$.
\choice
{$\text{IX}, \text{XV}, \text{XIX}, \text{XXI}$}
{\True $\text{XXI}, \text{XIX}, \text{XV}, \text{IX}$}
{$\text{XXI}, \text{XV}, \text{XIX}, \text{IX}$}
{$\text{XIX}, \text{XXI}, \text{XV}, \text{IX}$}
\loigiai{
Giá trị: $21 > 19 > 15 > 9$ nên thứ tự giảm dần là $\text{XXI}, \text{XIX}, \text{XV}, \text{IX}$.
\textbf{Đáp án của câu hỏi là B}
}
\end{ex}

\textbf{Phần 2. Trắc nghiệm đúng sai}
\setcounterpageref{ex}{0}

\begin{ex} %[3D1V49]
Xét tính đúng sai của các nhận định ôn tập chủ đề các số đến 10 000:
\choiceTF
{\True Số 10 000 gồm 10 nghìn hoặc 100 trăm.}
{Chữ số La Mã $\text{X}$ có giá trị bằng 15.}
{\True Khi so sánh hai số có 4 chữ số, ta luôn so sánh từ hàng nghìn trước.}
{Số 4550 làm tròn đến hàng trăm là 4500.}
\loigiai{
\begin{itemize}
    \item Ý a) Đúng: Cấu tạo số chuẩn.
    \item Ý b) Sai: Sai vì $\text{X} = 10$.
    \item Ý c) Đúng: Quy tắc so sánh.
    \item Ý d) Sai: Sai vì hàng chục là 5 nên làm tròn lên 4600.
\end{itemize}
\begin{itemize}
    \item a) \textbf{Đúng}.
    \item b) \textbf{Sai}.
    \item c) \textbf{Đúng}.
    \item d) \textbf{Sai}.
\end{itemize}
}
\end{ex}

\begin{ex} %[3D1V49]
Cho ba số: $A = 3450, B = 3540, C = 3045$:
\choiceTF
{\True Số bé nhất trong ba số là $C = 3045$.}
{\True Số lớn nhất trong ba số là $B = 3540$.}
{Làm tròn số $A$ đến hàng trăm ta được 3400.}
{\True Hiệu giữa số lớn nhất và số bé nhất là: $3540 - 3045 = 495$.}
\loigiai{
\begin{itemize}
    \item Ý a) Đúng: Hàng trăm bé nhất là 0.
    \item Ý b) Đúng: Hàng trăm lớn nhất là 5.
    \item Ý c) Sai: Sai vì hàng chục của $A$ là 5 nên làm tròn lên 3500.
    \item Ý d) Đúng: Phép tính chính xác.
\end{itemize}
\begin{itemize}
    \item a) \textbf{Đúng}.
    \item b) \textbf{Đúng}.
    \item c) \textbf{Sai}.
    \item d) \textbf{Đúng}.
\end{itemize}
}
\end{ex}

\textbf{Phần 3. Trả lời ngắn}
\setcounterpageref{ex}{0}

\begin{ex} %[3D1V49]
Tìm số liền sau của số 9999.
\shortans{10000}
\loigiai{
\textbf{Đáp số:} $10000$\\
$9999 + 1 = 10\,000$.
}
\end{ex}

\begin{ex} %[3D1V49]
Viết số tự nhiên tương ứng với chữ số La Mã $\text{XVI}$.
\shortans{16}
\loigiai{
\textbf{Đáp số:} $16$\\
$\text{XVI} = 16$.
}
\end{ex}

\begin{ex} %[3D1V49]
Làm tròn số 8456 đến hàng chục.
\shortans{8460}
\loigiai{
\textbf{Đáp số:} $8460$\\
Hàng đơn vị là 6 nên làm tròn thành 8460.
}
\end{ex}

\begin{ex} %[3D1V49]
Làm tròn số 6749 đến hàng trăm.
\shortans{6700}
\loigiai{
\textbf{Đáp số:} $6700$\\
Hàng chục là 4 nên làm tròn thành 6700.
}
\end{ex}

\begin{ex} %[3D1V49]
Tìm số bé nhất có bốn chữ số khác nhau có tổng các chữ số bằng 10.
\shortans{1027}
\loigiai{
\textbf{Đáp số:} $1027$\\
Nghìn 1, trăm 0, chục 2, đơn vị $10 - 3 = 7$. Số đó là 1027.
}
\end{ex}

\begin{ex} %[3D1V49]
Có bao nhiêu số tròn chục có bốn chữ số từ 2010 đến 2090?
\shortans{9}
\loigiai{
\textbf{Đáp số:} $9$\\
$(2090 - 2010) : 10 + 1 = 9$ số.
}
\end{ex}
